\documentclass[11pt,a4paper]{article}
\usepackage[margin=2.4cm]{geometry}
\usepackage{amsmath,amssymb}
\usepackage{booktabs}
\usepackage{xcolor}
\usepackage[colorlinks=true,linkcolor=blue,urlcolor=blue]{hyperref}

\newcommand{\ep}{\epsilon}
\newcommand{\zb}{\bar z}
\newcommand{\dd}{\mathrm{d}}
\newcommand{\Li}{\mathrm{Li}_2}
\newcommand{\F}{{}_2F_1}
\newcommand{\note}[1]{\par\smallskip\noindent\colorbox{gray!12}{\parbox{\dimexpr\linewidth-2\fboxsep}{\small #1}}\par\smallskip}

\title{The phase space of Drell--Yan real emission at fixed rapidity\\[2pt]
\large NLO, NNLO and N$^3$LO in one closed formula}
\author{DY\_P2B --- \texttt{phase\_space\_int/}}
\date{\today}

\begin{document}
\maketitle

\begin{abstract}
The simplest integral of every real-emission family is the one whose integrand is \emph{only} the cut
propagators: the phase space of $n$ massless partons plus the colour singlet $Q$, at fixed invariant mass $Q^2=m^2$
and fixed rapidity. For the triple real ($n=3$) it is \texttt{INT(topoN,1,1,1,1,1,0,\dots,0)}, the same in all
families. We compute it for any $n$ in $d=4-2\ep$, exactly in $z$, $u$ and $\ep$. The result
(eq.~\eqref{eq:main}) is a product of powers and one Gauss hypergeometric function, which equals 1 for $n=1$; for
$n=2$ it is the result of Anastasiou, Dixon, Melnikov, Petriello (hep-ph/0312266, eq.~(3.8)).
We explain every step, give the $\ep$ expansion and the soft limit, and compare with the NNLO computation in
\texttt{nnlo\_rr/integrals/IF1/\allowbreak boundaries.nb}. Everything is verified numerically by the scripts in
\texttt{checks/}.
\end{abstract}

\tableofcontents

\section*{Roadmap}
\addcontentsline{toc}{section}{Roadmap}
The computation has three ideas, one per section:
\begin{enumerate}
\item \textbf{Factorise} (Sect.~\ref{sec:strategy}). The $n$ massless partons only feel the total momentum
  $K=p_1+p_2-Q$ they share. Their phase space is a known function of $K^2$ alone. What remains is an integral
  over $Q$.
\item \textbf{Light-cone variables for $Q$} (Sect.~\ref{sec:lightcone}). With $Q=a\,p_1+b\,p_2+Q_T$ the
  rapidity constraint is \emph{linear} ($b=u\,a$) and the mass-shell constraint fixes $|Q_T|$. One integral, over
  $a$, is left.
\item \textbf{Recognise a hypergeometric function} (Sect.~\ref{sec:hyp}). The $a$ integrand is a product of three
  powers. Mapping the integration range to $[0,1]$ gives Euler's integral for $\F$.
\end{enumerate}
The answer is in Sect.~\ref{sec:result}. The rest of the note covers its use: $\ep$ expansion, soft limit, the
comparison with NNLO, and N$^3$LO.

\section{The integral and what each factor means}
\label{sec:def}
\subsection{Definition}
Incoming massless partons $p_1,p_2$, with $s=(p_1+p_2)^2=2\,p_1\!\cdot p_2$. Outgoing massless partons
$k_1,\dots,k_n$. The colour singlet carries what is left:
\[ Q=p_1+p_2-\sum_{i=1}^n k_i . \]
We study
\begin{equation}
  \boxed{\;\Phi_n(z,u)=\int\prod_{i=1}^n \dd^dk_i\,\delta^+(k_i^2)\;\;\delta^+(Q^2-m^2)\;\;
  \delta(Q\cdot p_1-u\,Q\cdot p_2)\;}
  \qquad z=\frac{m^2}{s},\qquad d=4-2\ep .
  \label{eq:def}
\end{equation}
Here $\delta^+(k^2)=\delta(k^2)\,\theta(k^0)$. Each factor has a meaning:
\begin{itemize}
\item $\delta^+(k_i^2)$: parton $i$ is real (on shell, positive energy). These are the cut propagators $1/k_i^2$.
\item $\delta^+(Q^2-m^2)$: the singlet has invariant mass $m$, i.e.\ we look at fixed $Q^2$. This is the cut of
  $1/(Q^2-m^2)$.
\item $\delta(Q\cdot p_1-u\,Q\cdot p_2)$: the singlet has fixed rapidity (see below). This is the fifth
  ``propagator'' of our families. It is linear in the momenta.
\end{itemize}
No factors of $2\pi$, $i$ or $e^{\ep\gamma_E}$ are included. Any other normalisation is an overall $\ep$-dependent
constant; Sect.~\ref{sec:nnlo} identifies the one used at NNLO.

\subsection{Why the last delta fixes the rapidity}
Work in the partonic centre-of-mass frame, with $p_1=\frac{\sqrt s}2(1,0,0,1)$ and
$p_2=\frac{\sqrt s}2(1,0,0,-1)$. Write $Q=(m_T\cosh Y,\,\vec Q_T,\,m_T\sinh Y)$ with $m_T^2=m^2+Q_T^2$. Then
\[
  Q\cdot p_1=\tfrac{\sqrt s}{2}\,(Q^0-Q^3)=\tfrac{\sqrt s}{2}\,m_T\,e^{-Y},\qquad
  Q\cdot p_2=\tfrac{\sqrt s}{2}\,(Q^0+Q^3)=\tfrac{\sqrt s}{2}\,m_T\,e^{Y},
\]
so $Q\cdot p_1-u\,Q\cdot p_2=0$ means $u=e^{-2Y}$: \emph{$u$ is the rapidity of $Q$}. Note that
$m_T$ drops out of the ratio.

\note{\textbf{The ``different delta''.} One can also fix the rapidity with $\delta(u-u_Q)$,
$u_Q=Q\cdot p_1/Q\cdot p_2$. The two are related by
$\delta(Q\cdot p_1-u\,Q\cdot p_2)=\delta(u-u_Q)/(Q\cdot p_2)$, since the derivative of the argument with
respect to $u$ is $Q\cdot p_2$. So \eqref{eq:def} is the $u$-distribution of the phase space weighted with
$1/Q\cdot p_2$. This is the factor \texttt{1/p2dotq} of the NNLO notebook.}

\subsection{Physical region and dimension}
The rapidity of an object of mass $m$ produced at energy $\sqrt s$ is bounded by $|Y|\le\frac12\ln\frac1z$. In
terms of $u$:
\[ z\le u\le 1/z . \]
We will re-derive this below, as the condition for the $a$ integral to be non-empty.

Dimension check: each $\int\dd^dk\,\delta^+(k^2)$ has mass dimension $d-2$, and each of the two other deltas
has $-2$. So $[\Phi_n]=n(d-2)-4=2(n-2-n\ep)$, and $\Phi_n\propto s^{\,n-2-n\ep}$ times a function of $z$ and $u$.

\subsection{Link with our integral families}
In \texttt{Code\_RRR} every momentum is incoming: $r_i=-k_i$, $q=-Q$. The five cut propagators are
\[
  r_1^2,\quad r_2^2,\quad r_3^2,\quad (p_1+p_2+r_1+r_2+r_3)^2-m^2=Q^2-m^2,\quad
  -q\cdot p_1+u\,q\cdot p_2=Q\cdot p_1-u\,Q\cdot p_2 .
\]
Hence the sector with only the cut lines,
\[\texttt{INT(topoN,1,1,1,1,1,0,0,0,0,0,0,0)},\]
is $\Phi_3$ for \emph{every} family, up to the normalisation convention. The pipeline uses $m^2=1$, so $s=1/z$.

\section{Strategy: factorise off the massless partons}
\label{sec:strategy}
Insert $1=\int\dd^dQ\,\delta^d\big(Q-p_1-p_2+\sum_ik_i\big)$ into \eqref{eq:def}. The partons now only appear in
\begin{equation}
  \Phi^{(0)}_n(K^2)\equiv\int\prod_{i=1}^n\dd^dk_i\,\delta^+(k_i^2)\;\delta^d\Big(K-\sum_ik_i\Big),
  \qquad K=p_1+p_2-Q,
\end{equation}
the phase space of $n$ massless particles with total momentum $K$. It is Lorentz invariant, so it depends only on
$K^2$. Therefore
\begin{equation}
  \Phi_n=\int\dd^dQ\;\delta^+(Q^2-m^2)\;\delta(Q\cdot p_1-u\,Q\cdot p_2)\;\Phi^{(0)}_n\big((p_1+p_2-Q)^2\big).
  \label{eq:fact}
\end{equation}
Two tasks remain: compute $\Phi^{(0)}_n(K^2)$ (Sect.~\ref{sec:massless}) and the $Q$ integral
(Sect.~\ref{sec:lightcone}).

\section{\texorpdfstring{The massless $n$-body phase space}{The massless n-body phase space}}
\label{sec:massless}
We need the solid angle of the unit sphere in $D$ dimensions, $\Omega_D=2\pi^{D/2}/\Gamma(D/2)$
(e.g.\ $\Omega_3=4\pi$).

\subsection{Two bodies, one of them massive}
Take one massless particle $k$ and one of mass $M$, in the rest frame of $K$ ($K=(\sqrt{K^2},\vec 0)$):
\[
  \Phi_2(K^2;M^2,0)=\int\dd^dk\,\delta^+(k^2)\,\delta^+\big((K-k)^2-M^2\big)
  =\int\frac{\dd^{d-1}k}{2|\vec k|}\,\delta\big(K^2-2\sqrt{K^2}\,|\vec k|-M^2\big).
\]
In the first step $\int\dd k^0\,\delta^+(k^2)=1/(2|\vec k|)$, and we used $(K-k)^2=K^2-2\sqrt{K^2}|\vec k|$ for
massless $k$. In spherical coordinates, $\dd^{d-1}k=\Omega_{d-1}|\vec k|^{d-2}\dd|\vec k|$. The delta fixes
$|\vec k|=E\equiv(K^2-M^2)/(2\sqrt{K^2})$, with Jacobian $1/(2\sqrt{K^2})$. This gives
\[
  \Phi_2(K^2;M^2,0)=\frac{\Omega_{d-1}}{4\sqrt{K^2}}\,E^{\,d-3}
  =\frac{\pi^{1-\ep}\,\Gamma(1-\ep)}{2\,\Gamma(2-2\ep)}\;(K^2)^{-1+\ep}\,(K^2-M^2)^{1-2\ep},
\]
where the duplication formula $\Gamma(1-\ep)\Gamma(\frac32-\ep)=2^{-1+2\ep}\sqrt\pi\,\Gamma(2-2\ep)$ was used to
write $\Omega_{3-2\ep}=2^{2-2\ep}\pi^{1-\ep}\Gamma(1-\ep)/\Gamma(2-2\ep)$. Two massless particles: set $M=0$,
$\Phi^{(0)}_2(K^2)=\frac{\pi^{1-\ep}\Gamma(1-\ep)}{2\Gamma(2-2\ep)}(K^2)^{-\ep}$.

\subsection{\texorpdfstring{Recursion to $n$ bodies}{Recursion to n bodies}}
Group the first $n-1$ partons into one ``particle'' of mass $M$: insert
$1=\int\dd M^2\int\dd^dP\,\delta^+(P^2-M^2)\,\delta^d(P-\sum_{i<n}k_i)$. Then
\[
  \Phi^{(0)}_n(K^2)=\int_0^{K^2}\dd M^2\;\Phi_2(K^2;M^2,0)\;\Phi^{(0)}_{n-1}(M^2).
\]
Assume $\Phi^{(0)}_{n-1}(M^2)=c_{n-1}(M^2)^{n-3-(n-2)\ep}$. With $M^2=K^2t$ the integral is a Beta function,
$\int_0^1\dd t\,t^{\,n-3-(n-2)\ep}(1-t)^{1-2\ep}=B\big((n-2)(1-\ep),2-2\ep\big)$, and
\begin{equation}
  \Phi^{(0)}_n(K^2)=c_n\,(K^2)^{\,n-2-(n-1)\ep},\qquad
  c_n=\frac{\pi^{(n-1)(1-\ep)}\,\Gamma(1-\ep)^n}{2^{n-1}\,\Gamma\big((n-1)(1-\ep)\big)\,\Gamma\big(n(1-\ep)\big)}
  \quad(n\ge2).
  \label{eq:cn}
\end{equation}
The cases needed:
\[
  \Phi^{(0)}_1=\delta^+(K^2),\qquad
  c_2=\frac{\pi^{1-\ep}\Gamma(1-\ep)}{2\Gamma(2-2\ep)},\qquad
  c_3=\frac{\pi^{2-2\ep}\Gamma(1-\ep)^3}{4\Gamma(2-2\ep)\Gamma(3-3\ep)} .
\]
The power of $K^2$ follows from dimensions: $n(d-2)-d=2(n-2-(n-1)\ep)$.

\section{\texorpdfstring{The integral over $Q$ in light-cone variables}{The integral over Q in light-cone variables}}
\label{sec:lightcone}
\subsection{Why light-cone variables}
Both remaining constraints are simple in the basis $\{p_1,p_2,\text{transverse}\}$:
\begin{equation}
  Q=a\,p_1+b\,p_2+Q_T,\qquad Q_T\cdot p_1=Q_T\cdot p_2=0,\qquad Q_T^2=-q_T^2\ \ (q_T\ge0).
\end{equation}
Since $p_1^2=p_2^2=0$:
\[
  Q\cdot p_1=\tfrac s2\,b,\qquad Q\cdot p_2=\tfrac s2\,a,\qquad Q^2=s\,a\,b-q_T^2 .
\]
The rapidity constraint is linear in $(a,b)$ and the mass-shell constraint involves $q_T$ only through $q_T^2$.

\subsection{The measure}
In the frame of Sect.~\ref{sec:def}, $Q^0=\frac{\sqrt s}2(a+b)$ and $Q^3=\frac{\sqrt s}2(a-b)$. Hence
\[\dd Q^0\dd Q^3=\frac s4\,\Big|\det\begin{pmatrix}1&1\\1&-1\end{pmatrix}\Big|\,\dd a\,\dd b=\frac s2\,\dd a\,\dd b,\]
and
\begin{equation}
  \dd^dQ=\frac s2\,\dd a\,\dd b\,\dd^{d-2}q_T .
\end{equation}

\subsection{Doing the two deltas}
\textbf{Rapidity.} $\delta(Q\cdot p_1-u\,Q\cdot p_2)=\delta\big(\frac s2(b-u\,a)\big)=\frac2s\,\delta(b-u\,a)$.
The $b$ integral sets $b=u\,a$, and the $\frac2s$ cancels the $\frac s2$ of the measure.

\textbf{Mass shell.} Now $Q^2=s\,u\,a^2-q_T^2$. In the $(d-2)$-dimensional Euclidean transverse space,
$\dd^{d-2}q_T=\frac{\Omega_{d-2}}2(q_T^2)^{\frac{d-4}2}\dd q_T^2$, so
\begin{equation}
  \int\dd^{d-2}q_T\;\delta\big(s u a^2-q_T^2-m^2\big)=\frac{\pi^{1-\ep}}{\Gamma(1-\ep)}\,\big(s u a^2-m^2\big)^{-\ep},
  \label{eq:qt}
\end{equation}
with $q_T^2=s u a^2-m^2$ fixed. At $\ep=0$ this is just $\pi$: in 4 dimensions the transverse ``sphere'' is a circle
and its length times the Jacobian is constant. The factor $(\dots)^{-\ep}$ is the $d$-dimensional remnant.

\textbf{The recoil.} With $b=ua$ and $q_T^2=sua^2-m^2$, the momentum of the partons
$K=(1-a)p_1+(1-ua)p_2-Q_T$ has
\[
  K^2=s(1-a)(1-ua)-q_T^2=s\big(1-a(1+u)\big)+m^2=s(1+u)\,(a_1-a),\qquad a_1\equiv\frac{1+z}{1+u}.
\]

\subsection{Integration limits and the physical region}
$a$ is bounded on both sides:
\begin{itemize}
\item $q_T^2\ge0$ requires $a\ge a_0\equiv\sqrt{z/u}$; $a>0$ comes from $Q^0>0$;
\item $K^2\ge0$ requires $a\le a_1$; then also $K^0=\frac{\sqrt s}2(2-a(1+u))>0$, since $a_1(1+u)=1+z<2$.
\end{itemize}
The range is non-empty iff $a_0\le a_1$. Multiplying by $(1+u)$, this is $w\ge0$, where
\begin{equation}
  w\equiv(1+u)(a_1-a_0)=1+z-(1+u)\sqrt{z/u}=\Big(1-\sqrt{\tfrac zu}\Big)\Big(1-\sqrt{zu}\Big).
\end{equation}
Both factors are $\ge0$ exactly for $z\le u\le1/z$ (both negative is impossible for $z<1$). This re-derives the
physical region; $w\to0$ at its edges, the maximal $|Y|$.

\subsection{Result: a one-dimensional integral}
Collect everything. Use $s u a^2-m^2=s u\,(a-a_0)(a+a_0)$ and $\Phi^{(0)}_n(K^2)=c_n[s(1+u)(a_1-a)]^{k}$ with
$k\equiv n-2-(n-1)\ep$:
\begin{equation}
  \Phi_n=\frac{\pi^{1-\ep}}{\Gamma(1-\ep)}\,c_n\;s^{\,n-2-n\ep}\,u^{-\ep}(1+u)^{k}
  \int_{a_0}^{a_1}\!\dd a\,(a-a_0)^{-\ep}\,(a+a_0)^{-\ep}\,(a_1-a)^{k}.
  \label{eq:oned}
\end{equation}
For $n=1$, $\Phi^{(0)}_1=\delta(K^2)$ instead puts $a=a_1$ and no integral is left (Sect.~\ref{sec:result}).

\section{The hypergeometric function}
\label{sec:hyp}
Map the range to $[0,1]$ with $a=a_0+(a_1-a_0)\,t$:
\[
  a-a_0=(a_1-a_0)\,t,\qquad a_1-a=(a_1-a_0)(1-t),\qquad a+a_0=2a_0\,(1+x\,t),\qquad x\equiv\frac{a_1-a_0}{2a_0}.
\]
The integral becomes
\[
  (a_1-a_0)^{\,n-1-n\ep}\,(2a_0)^{-\ep}\int_0^1\dd t\;t^{-\ep}(1-t)^{k}(1+x\,t)^{-\ep}.
\]
The two powers of $t$ and $1-t$ alone would give a Beta function. The extra $(1+xt)^{-\ep}$ turns it into Euler's
integral representation
\[
  \int_0^1\dd t\,t^{\beta-1}(1-t)^{\gamma-\beta-1}(1-wt)^{-\alpha}=B(\beta,\gamma-\beta)\;\F(\alpha,\beta;\gamma;w),
\]
here with $\alpha=\ep$, $\beta=1-\ep$, $\gamma=n-n\ep$, $w=-x$:
\begin{equation}
  \int_{a_0}^{a_1}\!\dd a\cdots=(a_1-a_0)^{\,n-1-n\ep}\,(2a_0)^{-\ep}\;
  B\big(1-\ep,(n-1)(1-\ep)\big)\;\F\big(\ep,1-\ep;n-n\ep;-x\big).
\end{equation}
The variable $x$ measures how far the integration range extends compared with its distance from the singular
point $a=-a_0$. Near threshold or at the edges of the rapidity range $x\to0$ and $\F\to1$.

\section{The result}
\label{sec:result}
Two identities make the result compact. First,
$u(a_1-a_0)(a_1+a_0)=u\,a_1^2-z=\frac{(u-z)(1-uz)}{(1+u)^2}$ (multiply out). Second,
$(2a_0)^{-\ep}=(a_1+a_0)^{-\ep}(1+x)^{\ep}$. Using them, together with
$(1+u)^k(a_1-a_0)^{(n-1)(1-\ep)}=w^{(n-1)(1-\ep)}/(1+u)$:
\begin{equation}
  \boxed{\;
  \Phi_n(z,u)=C_n\;s^{\,n-2-n\ep}\;\frac{w^{(n-1)(1-\ep)}}{1+u}\;
  \Big[\frac{(u-z)(1-uz)}{(1+u)^2}\Big]^{-\ep}\;G_n(\ep,x)\;}
  \label{eq:main}
\end{equation}
\begin{align}
  G_n(\ep,x)&=(1+x)^{\ep}\,\F\big(\ep,1-\ep;n-n\ep;-x\big), &
  C_n&=\frac{\pi^{n(1-\ep)}\,\Gamma(1-\ep)^n}{2^{n-1}\,\Gamma\big(n(1-\ep)\big)^2},\nonumber\\
  w&=\Big(1-\sqrt{\tfrac zu}\Big)\Big(1-\sqrt{zu}\Big), &
  x&=\frac{w\sqrt u}{2\sqrt z\,(1+u)} .
\end{align}
Here $C_n=\frac{\pi^{1-\ep}}{\Gamma(1-\ep)}\,c_n\,B(1-\ep,(n-1)(1-\ep))$.

\note{\textbf{Equivalent form, as in the literature.} By Pfaff's transformation (Sect.~\ref{sec:nnlo}),
$G_n=\F\big(\ep,(n-1)(1-\ep);n(1-\ep);\delta\big)$ with
$\delta=\frac{x}{1+x}=\frac{(\sqrt u-\sqrt z)(1-\sqrt{uz})}{(\sqrt u+\sqrt z)(1+\sqrt{uz})}\in[0,1)$. For $n=2$ this
is eq.~(3.8)--(3.11) of Anastasiou, Dixon, Melnikov, Petriello, hep-ph/0312266 (their $I[0]$, $K_0(\delta)$),
including the normalisation. For $n=3$ the formula is the generalisation to three partons.}

The structure of \eqref{eq:main}:
\begin{itemize}
\item $w^{(n-1)(1-\ep)}$: the ``available phase space'' for the partons. It vanishes at the kinematic boundary;
  each extra parton adds one power.
\item $[(u-z)(1-uz)/(1+u)^2]^{-\ep}$: from the transverse momentum of $Q$ in $d$ dimensions. At NLO this is the
  well-known $[(u-z)(1-uz)]^{-\ep}$.
\item $G_n$: the only non-trivial function, $G_n=1+O(\ep\,x)$.
\end{itemize}

\subsection{The three orders}
\begin{itemize}
\item \textbf{$n=1$ (NLO).} Here $\delta(K^2)=\delta\big(s(1+u)(a_1-a)\big)$ sets $a=a_1$. Directly from
  \eqref{eq:qt}:
  \[\Phi_1=\frac{\pi^{1-\ep}}{\Gamma(1-\ep)}\;\frac{s^{-1-\ep}}{1+u}\;\Big[\frac{(u-z)(1-uz)}{(1+u)^2}\Big]^{-\ep}.\]
  This is \eqref{eq:main} with $C_1=\pi^{1-\ep}/\Gamma(1-\ep)$, $w^0=1$ and $G_1=1$. Consistently,
  $\F(\ep,1-\ep;1-\ep;-x)=(1+x)^{-\ep}$, so $G_1=1$ also from the formula. At NLO the deltas fix everything; nothing
  is left to integrate.
\item \textbf{$n=2$ (NNLO).}
  $\Phi_2=\frac{\pi^{2-2\ep}\Gamma(1-\ep)^2}{2\Gamma(2-2\ep)^2}\,s^{-2\ep}\,
  \frac{w^{1-\ep}}{1+u}\big[\tfrac{(u-z)(1-uz)}{(1+u)^2}\big]^{-\ep}G_2(\ep,x)$.
\item \textbf{$n=3$ (N$^3$LO).}
  \[\Phi_3=\frac{\pi^{3-3\ep}\Gamma(1-\ep)^3}{4\,\Gamma(3-3\ep)^2}\;s^{1-3\ep}\;
  \frac{w^{2-2\ep}}{1+u}\Big[\frac{(u-z)(1-uz)}{(1+u)^2}\Big]^{-\ep}(1+x)^{\ep}\,\F(\ep,1-\ep;3-3\ep;-x).\]
\end{itemize}

\subsection{Sanity checks you can do by eye}
\begin{itemize}
\item \textbf{$\ep=0$:} $G_n=1$ and $\Phi_n=C_n\,s^{n-2}w^{n-1}/(1+u)$, e.g.\ $\Phi_3=\frac{\pi^3}{16}\,s\,\frac{w^2}{1+u}$.
\item \textbf{Dimension:} $s^{\,n-2-n\ep}$, as counted in Sect.~\ref{sec:def}.
\item \textbf{Edges of the rapidity range:} $w\to0$ at $u\to z$ and $u\to1/z$; the phase space closes.
\item \textbf{$p_1\leftrightarrow p_2$:} exchanging the beams sends $Y\to-Y$, i.e.\ $u\to1/u$. From the delta,
  $\delta(Q\cdot p_1-\tfrac1u Q\cdot p_2)=u\,\delta(Q\cdot p_2-u\,Q\cdot p_1)$, so one needs
  $\Phi_n(z,1/u)=u\,\Phi_n(z,u)$. In \eqref{eq:main}, $w$, $x$ and $(u-z)(1-uz)/(1+u)^2$ are invariant under
  $u\to1/u$, and $\frac1{1+1/u}=\frac u{1+u}$. \checkmark
\end{itemize}

\section{\texorpdfstring{The $\ep$ expansion}{The epsilon expansion}}
\label{sec:eps}
All factors of \eqref{eq:main} except $G_n$ are powers: they expand into logarithms of $w$, $1+u$,
$(u-z)(1-uz)$ and constants. For $G_n$, expand the series of $\F$ or the Euler integral. With $L=\ln(1+x)$:
\begin{align}
G_2&=1+\ep\Big(1-\frac Lx\Big)+\ep^2\,\frac{6x-2L+xL^2+2(2+x)\Li(-x)}{2x}+O(\ep^3),\\
G_3&=1+\ep\,\frac{x(2+3x)-2(1+2x)L}{2x^2}\nonumber\\
 &\quad+\ep^2\,\frac{x(7+11x)-(1+2x)L+(1+2x+2x^2)L^2+2(3+6x+2x^2)\Li(-x)}{2x^2}+O(\ep^3).
\end{align}
How the $O(\ep)$ term arises: $(\ep)_m\approx\ep\,(m-1)!$ for small $\ep$, so
$\F(\ep,1-\ep;n;-x)\approx1+\ep\sum_{m\ge1}\frac{(m-1)!}{(n)_m}(-x)^m$. For $n=2$ the sum is
$\sum_m\frac{(-x)^m}{m(m+1)}=1-\frac{(1+x)L}{x}$; adding $\ep L$ from $(1+x)^\ep$ gives $1-L/x$. These
coefficients are regular at $x=0$ (e.g.\ $1-L/x=x/2+O(x^2)$), so $G_n-1$ is small near threshold.
Higher orders: expand the Euler integral; this gives harmonic polylogarithms of $x$ (HypExp or PolyLogTools).

\section{NNLO variables and the soft limit}
\label{sec:soft}
\subsection{The variables of the NNLO computation}
The NNLO computation uses $\zb=1-z$ and $y\in[0,1]$, with
\[ u=\frac{y+z-yz}{1-y+yz}\qquad(y=0\Rightarrow u=z,\quad y=1\Rightarrow u=1/z). \]
So $y$ maps the rapidity range onto $[0,1]$. A short computation gives the exact relations
\begin{gather}
  1+u=\frac{1+z}{1-y\zb},\qquad u-z=\frac{y(1-z^2)}{1-y\zb},\qquad 1-uz=\frac{(1-y)(1-z^2)}{1-y\zb},\nonumber\\
  \frac{(u-z)(1-uz)}{(1+u)^2}=y(1-y)\,\zb^{\,2}.
\end{gather}
So the bracket in \eqref{eq:main} is rational in $(y,z)$; the square roots only sit in $w$ and $x$. These are
the roots the NNLO notebook handles with $r=\sqrt u$, $t=\sqrt z$.

\subsection{Soft limit}
$\zb\to0$ is the threshold: no energy is left for radiation. At fixed $y$ one finds
$\sqrt{z/u}\simeq1-y\zb$ and $\sqrt{zu}\simeq1-(1-y)\zb$, hence
\[
  w=y(1-y)\zb^2\big(1+\tfrac{\zb}2+O(\zb^2)\big),\qquad x=\tfrac14\,y(1-y)\zb^2+O(\zb^3),\qquad
  G_n=1+O(\ep\,\zb^2),\qquad 1+u\to2 .
\]
With $m^2=1$, $s^{\dots}=z^{-\dots}\to1$, and the leading power is
\begin{equation}
  \Phi_n^{\rm LP}=\frac{C_n}{2}\,\big(y(1-y)\big)^{\,n-1-n\ep}\,\zb^{\,2(n-1)-2n\ep}.
  \label{eq:lp}
\end{equation}
The exponents count the phase space: each parton contributes $\zb^{2-2\ep}$ (energy and angles scale with $\zb$),
and $\big(y(1-y)\zb^2\big)^{-\ep}$ comes from the transverse momentum of $Q$.

\section{Comparison with the NNLO computation}
\label{sec:nnlo}
File: \texttt{nnlo\_rr/integrals/IF1/boundaries.nb}, section ``phase space''. There the phase-space master
$j[\mathrm{IF1},1,1,1,1,0,0,0]$ is integrated over angular variables $\lambda_1,\lambda_2,\lambda_4$ and expanded
at leading power in $\zb$. The exported result is the boundary condition of the canonical differential equations.

\subsection{(a) The exported soft master is right}
The file \texttt{SoftMasters/m1\_1111000.m} contains
\texttt{P[bz,2-4eps]\,P2[(1-y)y,1-2eps]} times a series in $\ep$ up to $\ep^6$. The prefactor is exactly the
structure of \eqref{eq:lp} for $n=2$. Its coefficients agree \emph{exactly, through $\ep^6$} with
\begin{equation}
  \frac{e^{2\ep\gamma_E}\,\Gamma(1-\ep)^2}{4\,\Gamma(2-2\ep)^2}
  =\frac{e^{2\ep\gamma_E}}{\pi^{2-2\ep}}\times\frac{C_2}{2}
  =\frac14+\ep+\ep^2\Big(3-\frac{\pi^2}8\Big)+\dots
\end{equation}
So the NNLO boundary is the leading power of the exact $\Phi_2$. This also fixes the normalisation used at NNLO:
$\mathcal N_2=e^{2\ep\gamma_E}/\pi^{2-2\ep}$ times \eqref{eq:def}, i.e.\ a factor $e^{\ep\gamma_E}/\pi^{1-\ep}$ per
real parton.

\subsection{(b) The closed formula in the comment: ADMP eq.~(3.8), exact}
The notebook also contains, in a comment (``result to print''), the closed formula
\[
  \frac{\Omega_{d-2}\,\Omega_{d-1}}{2^d}\;z^{2\ep}\,
  \frac{\big((u-z)(1-uz)\big)^{1-2\ep}}{(1+u)^{1-2\ep}\,\big(\sqrt u(\sqrt u+\sqrt z)(1+\sqrt{uz})\big)^{1-\ep}}
  \;K[0]\qquad(\times\,e^{2\ep\gamma_E}\pi^{2-d}).
\]
This is eq.~(3.8) of Anastasiou, Dixon, Melnikov, Petriello, \emph{High-precision QCD at hadron colliders:
electroweak gauge boson rapidity distributions at NNLO}, hep-ph/0312266 (ADMP), for $\nu=0$, with
\begin{gather}
  K[0]=K_0(\delta)=\int_0^1\dd\chi\,\big[\chi(1-\chi)(1-\chi\delta)\big]^{-\ep}
  =B(1-\ep,1-\ep)\;\F(\ep,1-\ep;2-2\ep;\delta),\nonumber\\
  \delta=\frac{(\sqrt u-\sqrt z)(1-\sqrt{uz})}{(\sqrt u+\sqrt z)(1+\sqrt{uz})} .
  \label{eq:admp}
\end{gather}
$K_0$ is \emph{a function of $z$ and $u$}, not a constant. With it the formula is exact, and it coincides with ours:
\begin{itemize}
\item \textbf{Same variable.} Our $x=(a_1-a_0)/(2a_0)$ and ADMP's $\delta$ are related by
  $\delta=\frac{x}{1+x}=\frac{a_1-a_0}{a_1+a_0}$ (multiply out $a_1\pm a_0$; $(\sqrt u+\sqrt z)(1+\sqrt{uz})
  =\sqrt u(1+z)+\sqrt z(1+u)$).
\item \textbf{Same function.} Pfaff's transformation $\F(a,b;c;w)=(1-w)^{-a}\F\big(a,c-b;c;\frac{w}{w-1}\big)$ with
  $w=-x$ gives
  \begin{equation}
    G_n=(1+x)^{\ep}\F(\ep,1-\ep;n-n\ep;-x)=\F\big(\ep,(n-1)(1-\ep);n(1-\ep);\delta\big),
    \label{eq:pfaff}
  \end{equation}
  so for $n=2$: $G_2=\F(\ep,1-\ep;2-2\ep;\delta)=K_0(\delta)/B(1-\ep,1-\ep)$.
\item \textbf{Same prefactor.} As shown in (a'), the power structure equals ours:
  $\frac{(u-z)(1-uz)}{\sqrt u(\sqrt u+\sqrt z)(1+\sqrt{uz})}=w$.
\item \textbf{Same normalisation.} $\Omega_{d-2}\Omega_{d-1}/2^d=\pi^{2-2\ep}/(2\Gamma(2-2\ep))=C_2/B(1-\ep,1-\ep)$
  (duplication formula), and ADMP's measure (their eq.~(3.6)) is our \eqref{eq:def}. Numerically
  $I[0]_{\rm ADMP}=\Phi_2$ to 30 digits.
\item \textbf{Same expansion.} ADMP's eq.~(3.12), $K_0=1+\ep\big(3-\frac{(\delta-1)\ln(1-\delta)}{\delta}\big)+\dots$,
  is $B(1-\ep,1-\ep)\,G_2$ with our $G_2=1+\ep(1-\ln(1+x)/x)+\dots$ ($B=1+2\ep+\dots$).
\end{itemize}

\paragraph{(a') The power structure.} Write $(u-z)=(\sqrt u-\sqrt z)(\sqrt u+\sqrt z)$ and
$(1-uz)=(1-\sqrt{uz})(1+\sqrt{uz})$; then
$\frac{(u-z)(1-uz)}{\sqrt u(\sqrt u+\sqrt z)(1+\sqrt{uz})}=\frac{(\sqrt u-\sqrt z)(1-\sqrt{uz})}{\sqrt u}=w$, and the
comment formula without $K[0]$ is
$F_c\equiv\frac{w^{1-\ep}}{1+u}\big[\frac{(u-z)(1-uz)}{(1+u)^2}\big]^{-\ep}$, i.e.\ \eqref{eq:main} for $n=2$
without $G_2$.

\paragraph{What changes if $K[0]$ is taken as a constant.} Only then would the formula be approximate: $K_0(\delta)$
varies with $(z,u)$ through $\delta$, and $\delta\to0$ only at threshold ($\delta\simeq x\simeq\frac14y(1-y)\zb^2$).
The table shows $K_0/B=G_2$ at $\ep=0.3$ (the middle column is a direct numerical integration of \eqref{eq:oned}
divided by $C_2s^{-2\ep}F_c$):

\begin{center}
\begin{tabular}{cccccc}
\toprule
$\zb$ & $y$ & $x$ & $\Phi_2/(C_2 s^{-2\ep}F_c)$ (numerical) & $G_2(\ep,x)$ & $1+\ep x/2$\\
\midrule
$1/2$   & $1/5$ & $1.96\cdot10^{-2}$ & $1.0029115462$ & $1.0029115462$ & $1.00294$\\
$1/2$   & $1/2$ & $3.03\cdot10^{-2}$ & $1.0044765289$ & $1.0044765289$ & $1.00455$\\
$1/10$  & $1/5$ & $4.44\cdot10^{-4}$ & $1.0000666211$ & $1.0000666211$ & $1.000067$\\
$1/10$  & $1/2$ & $6.94\cdot10^{-4}$ & $1.0001040555$ & $1.0001040555$ & $1.000104$\\
$1/100$ & $1/5$ & $4.04\cdot10^{-6}$ & $1.0000006061$ & $1.0000006061$ & $1.00000061$\\
$1/100$ & $1/2$ & $6.31\cdot10^{-6}$ & $1.0000009470$ & $1.0000009470$ & $1.00000095$\\
\bottomrule
\end{tabular}
\end{center}
At leading power $K_0\to B(1-\ep,1-\ep)$, a constant: that is what the notebook evaluates
(\texttt{Series[\dots,\char`\{bz,0,0\char`\}]}) and exports, and (a) confirms it.

\subsection{(c) The notebook integrand, with both deltas}
We also integrated the notebook's own \texttt{INTEGRAND} (copied verbatim, with the $(y,z)$-dependent part of
\texttt{Norm}) over $\lambda_1,\lambda_2,\lambda_4$. The $\lambda_4$ and $\lambda_2$ integrals factorise into Beta
functions; $\lambda_1$ is integrated numerically. Two objects:
\[
  A=\int\texttt{INTEGRAND}\quad(\text{delta }\delta(u-u_Q)),\qquad
  B=\int\texttt{INTEGRAND}/Q\!\cdot\!p_2\quad(\text{the master}).
\]
Ratios over five points $(z,u)$; a ratio that is the same at every point means the two functions coincide up to an
$\ep$-dependent constant:

\begin{center}
\begin{tabular}{lcc}
\toprule
ratio & $\ep=0$ & $\ep=0.3$\\
\midrule
$A\,/\,$\eqref{eq:oned} with extra factor $Q\cdot p_2=\frac{as}2$ & $2\pi$ everywhere & $36.0874\ldots$ everywhere\\
$B\,/\,$\eqref{eq:oned} (the master) & $2\pi$ everywhere & $36.0874\ldots$ everywhere\\
$A\,/\,F_c$ & $1.34\ \dots\ 2.91$ & $14.7\ \dots\ 31.7$\\
$B\,/\,F_c$ & $2\pi$ everywhere & $68.5\ \dots\ 72.7$\\
\bottomrule
\end{tabular}
\end{center}
The notebook's parametrisation is \emph{exact for all $z$}, for both deltas, and the factor $1/Q\cdot p_2$ pulled out
of the delta is handled correctly. $B/F_c$ is not constant at $\ep\ne0$ only because $F_c$ lacks the function
$K_0(\delta)$; $B/(F_cK_0)$ is constant.

\section{\texorpdfstring{N$^3$LO: what to use}{N3LO: what to use}}
\label{sec:n3lo}
\begin{itemize}
\item \textbf{The master.} \texttt{INT(topoN,1,1,1,1,1,0,\dots,0)}$\;=\Phi_3$ of \eqref{eq:main}, exact. In the
  normalisation analogous to NNLO, $\mathcal N_3=e^{3\ep\gamma_E}/\pi^{3-3\ep}$ (an assumption, to be matched to the
  convention of the reduction and the other masters):
  \[ \mathcal N_3\,\Phi_3=\frac{e^{3\ep\gamma_E}\Gamma(1-\ep)^3}{4\,\Gamma(3-3\ep)^2}\;s^{1-3\ep}\;
     \frac{w^{2-2\ep}}{1+u}\,\big[y(1-y)\zb^2\big]^{-\ep}\,G_3(\ep,x).\]
\item \textbf{Soft boundary} (if the NNLO strategy is repeated for the other masters):
  \[\mathcal N_3\Phi_3^{\rm LP}=\big(y(1-y)\big)^{2-3\ep}\zb^{4-6\ep}\times
  \frac{e^{3\ep\gamma_E}\Gamma(1-\ep)^3}{8\Gamma(3-3\ep)^2},\]
  \begin{align*}\frac{e^{3\ep\gamma_E}\Gamma(1-\ep)^3}{8\Gamma(3-3\ep)^2}&=\frac1{32}+\frac{9\ep}{32}
   +\frac{207-5\pi^2}{128}\ep^2+\frac{972-45\pi^2-68\zeta_3}{128}\ep^3\\
   &\quad+\frac{488430-31050\pi^2+163\pi^4-73440\zeta_3}{15360}\ep^4+\dots\end{align*}
  For \emph{this} master no boundary is needed, since the exact formula is known. It can also serve as an exact
  check of the differential-equation solution at any $z$.
\end{itemize}

\section{Checks}
\label{sec:checks}
\textbf{\texttt{checks/checks.wl}} (Mathematica; output in \texttt{checks.log}):
\begin{enumerate}
\item the one-dimensional integral \eqref{eq:oned} against the $\F$ form and against \eqref{eq:main}, for $n=2,3,4$
  and five points $(\ep,z,u)$ with $\ep$ of both signs: agreement to $10^{-28}$;
\item $n=1$: $G_1=1$ (symbolic);
\item the identities of the NNLO variables (symbolic);
\item the leading power of $w$ and $x$;
\item the $\ep$ and $\ep^2$ coefficients of $G_2$ and $G_3$: from the series of $\F$ and the Euler integral, checked
  by finite differences; the simplified forms above are identical to the Euler-integral results;
\item the NNLO soft master file against $e^{2\ep\gamma_E}\Gamma(1-\ep)^2/(4\Gamma(2-2\ep)^2)$: identical through $\ep^6$;
\item the comment formula of \texttt{boundaries.nb}: symbolic identity with $w^{1-\ep}/(1+u)[\dots]^{-\ep}$, and the
  table of Sect.~\ref{sec:nnlo};
\item the N$^3$LO leading-power constant;
\item the $p_1\leftrightarrow p_2$ symmetry $\Phi_n(z,1/u)=u\,\Phi_n(z,u)$.
\end{enumerate}
\textbf{\texttt{checks/admp.wl}} (output in \texttt{admp.log}): ADMP's $I[0]$, eq.~(3.8)--(3.11), against
$\Phi_2$ in absolute normalisation (30 digits); the Pfaff form \eqref{eq:pfaff} for $n=2,3$; $\delta=x/(1+x)$
(symbolic); the $O(\ep)$ term of $K_0$.

\textbf{\texttt{checks/notebook\_integrand.wl}} (output in \texttt{notebook\_integrand.log}): the notebook's
\texttt{INTEGRAND}, with and without $1/Q\cdot p_2$, against the light-cone results and $F_c$ (table of (c)).
\textbf{\texttt{checks/rest\_frame.py}} (plain Python; output in \texttt{rest\_frame.log}): an independent
derivation in the partonic rest frame instead of light-cone variables. $P\to Q+K$ two-body, $K\to n$ massless as in
Sect.~\ref{sec:massless}, and the angle of $Q$ fixed by the rapidity delta. Compared with \eqref{eq:main} in absolute
normalisation for $n=2,3$, $\ep=0,\,0.3,\,-0.2$ and three $(s,z,u)$ points: agreement to $10^{-9}$--$10^{-14}$ (the
accuracy of the quadrature).

\end{document}
